% year: תשפ"ד
% type: מבחן
% semester: א
% moed: מיוחד
% number: 4
% points: 14
% categories: ivt, mvt, func-analysis
% source: exams/מבחנים/תשפד/מועד ג תשפד.pdf

\begin{question}
הוכיחו כי למשוואה הבאה יש פתרון ממשי חיובי יחיד:
\[ 2x^3+\arctan x+\sqrt{x}=7 \]
\end{question}

\begin{hint}
קיום: משפט ערך הביניים עבור $g(x)=2x^3+\arctan x+\sqrt x-7$ בקטע מתאים. יחידות: הראו ש-$g$ עולה ממש.
\end{hint}

\begin{solution}
נגדיר $g(x)=2x^3+\arctan x+\sqrt x-7$ על $[0,\infty)$. $g$ רציפה שם כסכום של פונקציות רציפות, וגזירה ב-$(0,\infty)$.

\textbf{קיום.} $g(0)=-7<0$, ו-
\[ g(2)=16+\arctan2+\sqrt2-7=9+\arctan 2+\sqrt2>0. \]
לפי משפט ערך הביניים קיימת $x_0\in(0,2)$ עם $g(x_0)=0$, כלומר פתרון חיובי של המשוואה.

\textbf{יחידות.} לכל $x>0$:
\[ g'(x)=6x^2+\frac1{1+x^2}+\frac1{2\sqrt x}>0. \]
לכן $g$ עולה ממש ב-$(0,\infty)$ (מסקנה ממשפט לגרנז': אם $0<x_1<x_2$ אז $g(x_2)-g(x_1)=g'(c)(x_2-x_1)>0$). פונקציה עולה ממש מקבלת כל ערך לכל היותר פעם אחת, ולכן יש לה לכל היותר אפס אחד ב-$(0,\infty)$.

(לחלופין: אם היו שני פתרונות חיוביים $x_1<x_2$, אז לפי משפט רול הייתה $c\in(x_1,x_2)$ עם $g'(c)=0$, בסתירה ל-$g'>0$.)

\textbf{מסקנה:} למשוואה יש פתרון ממשי חיובי יחיד (מספרית $x_0\approx1.3486$).
\end{solution}
